% =====================================================================
% IEEEtran 会议论文最小可编译骨架 (bare_conf.tex 风格)
% 编译: latexmk -pdf ieee_bare_conf.tex      (pdfLaTeX, 引擎默认即可)
% 参考文献: 把 refs.bib 放同目录, 取消 \bibliography 注释后再编译
% 双盲: 去掉 \author 内真实信息, 或用占位
% =====================================================================
\documentclass[conference]{IEEEtran}
% \documentclass[conference,compsoc]{IEEEtran}   % IEEE Computer Society 变体
\IEEEoverridecommandlockouts                      % 允许 \thanks 等

\usepackage{cite}                 % IEEE 文献引用压缩 [1]-[3]
\usepackage{amsmath,amssymb,amsfonts}
\usepackage{algorithmic}
\usepackage{graphicx}
\usepackage{textcomp}
\usepackage{xcolor}
\usepackage[hidelinks]{hyperref}  % 放靠后; cleveref 若用须更靠后
\usepackage[capitalize]{cleveref}

\begin{document}

\title{Paper Title Here\\
{\footnotesize \textsuperscript{*}Optional subtitle / funding note}}

\author{\IEEEauthorblockN{First Author}
\IEEEauthorblockA{\textit{dept. name of org.} \\
\textit{name of org.}\\
City, Country \\
email@domain.com}
\and
\IEEEauthorblockN{Second Author}
\IEEEauthorblockA{\textit{dept. name} \\
\textit{name of org.}\\
City, Country \\
email@domain.com}}

\maketitle

\begin{abstract}
This is the abstract. State problem, method, key results and contributions in one paragraph.
\end{abstract}

\begin{IEEEkeywords}
keyword1, keyword2, keyword3
\end{IEEEkeywords}

\section{Introduction}
\IEEEPARstart{T}{his} is a drop-cap example. Citation example \cite{ref1}.
Cross-references: see \cref{fig:arch} and \cref{eq:main}.

\section{Method}
\begin{equation}\label{eq:main}
  y = f(x) + \epsilon
\end{equation}

\begin{figure}[t]
  \centering
  % \includegraphics[width=\linewidth]{fig1.pdf}
  \rule{0.6\linewidth}{3cm}  % 占位框, 有图时替换为 \includegraphics
  \caption{System architecture.}
  \label{fig:arch}
\end{figure}

% 宽图跨双栏放页顶:
% \begin{figure*}[t] \centering ... \caption{} \label{} \end{figure*}

\section{Conclusion}
Summary here.

% --- 参考文献 (二选一) ---------------------------------------------
% 方式 A: 手写 thebibliography (无需 bibtex, 适合演示)
\begin{thebibliography}{00}
\bibitem{ref1} A. Author, ``Title of paper,'' \textit{Journal}, vol.~1, no.~1, pp.~1--10, 2024.
\end{thebibliography}
% 方式 B: bibtex (推荐). 注释掉上面的 thebibliography, 启用下面两行:
% \bibliographystyle{IEEEtran}
% \bibliography{refs}

\end{document}
